\section{Introduction}
\label{sec:introduction}

Large language models (LLMs) make it possible to construct agents that
interpret natural-language situations, maintain plans, and generate socially
legible action.  This has renewed interest in generative agents and in using
LLMs to simulate individual or collective behavior
\citep{park2023generative,park2024grounded,argyle2023outofone}.  A convincing
conversation, however, is not yet an artificial society.  Society-level
research requires agents to persist across simulated time, act from locations
and relationships in a shared world, encounter consequences of earlier
actions, and leave evidence from which a researcher can reconstruct what
happened.  It also requires mechanisms outside the language model---time,
space, exposure, institutions, resource constraints, and interventions---to
remain explicit enough to inspect and vary.

These requirements create a systems problem.  Persistence cannot reside only
in a model context window: profiles, state, episodic experience, relationships,
and schedules must survive across ticks and runs.  Situatedness cannot be
reduced to scenery in a prompt: movement, co-location, local conditions, and
social ties must constrain later computation.  Experimental control requires
scenario configuration and intervention points that are distinct from the
agent's generated narrative.  Finally, scientific use requires provenance
below a polished transcript or aggregate plot.  A researcher should be able to
ask which mechanism produced a state change, whether every treatment
completed, and which model and configuration generated an archived trace.

GAWorld addresses this problem as research infrastructure for building
persistent and situated LLM agent societies.  It connects a continuing agent
lifecycle---profile, perception, memory and recall, planning and action,
reflection, state update, and storage---to an urban and societal substrate.
That substrate includes a city graph and local physical conditions, social
relationships and life events, economic state, work, information exposure,
and scheduled interventions.  A simulation clock and event boundaries connect
individual cognition to day-level and society-level evolution.  The resulting
unit of study is not an isolated model response but a configured run whose
local decisions, shared events, persistent state, and derived comparisons can
be inspected together.

The system is also an account of incremental research-software
modularization.  GAWorld began with tightly connected domain paths and is being
migrated toward a small execution kernel, configurable cognition stages, an
event bus, and lifecycle-managed plugins.  The current implementation contains
working examples of these boundaries while retaining direct calls, shared
dictionaries, subsystem-owned files, and hook-based compatibility paths.  We
document that mixed state rather than presenting an aspirational architecture
as complete.  This distinction is important for reuse and for science: a
mechanism that is configurable is not necessarily isolated, and an extension
that writes a metric is not necessarily on the causal path of an agent action.

This paper makes four contributions:

\begin{enumerate}[leftmargin=*]
  \item \textbf{A persistent and situated agent lifecycle.}  We describe how
  GAWorld joins relatively stable profiles, mutable structured state,
  file-backed memory, intentions, habits, relationships, interests, and skills
  to a configurable per-tick cognition pipeline and a shared simulation clock.
  The design makes continuity and context inspectable without treating them as
  evidence of psychological fidelity.

  \item \textbf{A composable urban-society substrate.}  We show how city,
  environment, social-network, economic, work, policy, and information
  mechanisms contribute to a common run.  Explicit coupling and provenance
  support experiments that cross individual and aggregate levels while making
  clear where calibration and mechanism validation remain necessary.

  \item \textbf{An extensible execution and research architecture.}  We
  present the configured stage pipeline, simulation context, event bus, plugin
  registry, controller and recorder skeletons, compatibility paths, and the
  command-line, dashboard, interview, replay, and relay interfaces built around
  them.  A migration-state audit separates wired services, partial
  integrations, and planned boundaries.

  \item \textbf{An inspectable experimentation workflow with bounded evidence.}
  We connect scenario configuration, matched-run comparison, structured
  artifacts, and GAWorld-Bench to an evidence ledger.  Existing memory,
  policy, information, and economy--wellbeing archives demonstrate what the
  system can record and expose; incomplete and diagnostic artifacts are
  reported alongside them rather than converted into claims of social
  validity.
\end{enumerate}

GAWorld is not presented as a digital twin, a representative synthetic
population, or a forecasting system.  Its current artifacts do not establish
that generated agents reproduce particular people, that aggregate trajectories
match a real city, or that a simulated intervention estimates a real policy
effect.  We instead treat the system as a scientific instrument whose outputs
must be evaluated at multiple levels and whose reproducibility is a prerequisite
for substantive interpretation.  This framing places software architecture,
experimental design, and validity boundaries in the same account.

The remainder of the paper proceeds from this distinction.  We first state
the design goals and position GAWorld relative to generative-agent,
language-agent, social-evaluation, and agent-based-modeling research.  We then
describe the agent lifecycle, execution architecture, memory and cognition,
and societal substrate.  Subsequent sections cover experiment and deployment
interfaces, audit the existing capability cases, and state the validation,
ethical, and development boundaries that govern intended use.
